\minitopic{依赖是复杂知识的常态}没有人能够独立验证药品研发、气候模型、桥梁设计和历史档案的全部环节。若把认识自主理解成不依赖他人，现代知识几乎都将失去资格。更合理的自主概念要求主体能够识别自己依赖谁、依赖什么能力、在何种条件下撤回信任。它不要求掌握专家的一阶技术证据，却要求拥有评估专家的元证据，并知道何时引入第二来源、何时等待制度复核。

\minitopic{初始信任与后来校准形成时间结构}证言非还原论认为，理解一项看似真诚断言即可获得可被击败的初步接受权利；还原论要求该权利最终由知觉、记忆和归纳支持。儿童学习显示，彻底先验审查不可能成为信任起点；但成年人持续信任明显失准来源又不能只诉诸初始权利。较完整的模型把二者放入发展过程：主体以有限初始信任进入语言共同体，通过纠错、交叉核验和来源记录逐步校准；当具体反证出现，初始权利收缩。争论便从“一律信或一律疑”转向信任怎样获得、维持和丧失。

\minitopic{来源数量必须按独立性折算}十篇报道若都来自同一匿名帖子，只是一个来源的十次显示；两家实验室若共享同一样本污染，也不能简单相乘证据。来源谱系要记录原始接触、采集、编辑、转述和平台分发，标出共享节点。独立性不是绝对隔离：研究团队可以共享背景理论，却使用独立测量；记者可以访问同一现场，却采访不同目击者。判断重点是给定核心假设后，一项来源的错误是否使另一项更容易同时错误。

\minitopic{专家评价不能由立场一致替代}非专家常无法直接判断复杂论证，只能考察领域训练、方法透明、同业记录、利益披露与纠错表现。这些元证据也可能冲突：一位专家履历优秀却拒绝公开数据，另一位缺少名望但方法可复现。评价应与任务匹配，不能用一般声望覆盖细分领域，也不能把同意自己的结论当作可靠标志。同侪分歧进一步要求区分一阶证据与“另一位能力相当者不同意”这一高阶证据；调整幅度取决于对方是否真正同侪、信息是否独立以及双方过去校准表现。

\begin{argumentbox}
社会证据档案至少包含：每项主张的直接来源；来源之间的共享节点；说话者在目标任务上的能力与利益；异议是否被记录和回应；结论由谁以何种程序正式接受；出现错误后谁能触发复核。把“很多人都这么说”替换成这六栏，群体判断才具有可审计结构。
\end{argumentbox}

\minitopic{群体态度不能等同成员多数}委员会可能依正式程序接受一项报告，即使部分成员私人反对；也可能人人私下相信，却没有任何群体决议。群体知识理论需要说明正式接受、共同承诺、信息聚合和功能结构中的哪一项构成群体态度。多数规则只是众多机制之一，且在信息相关、群体极化或权力不平等时可能放大错误。一个群体是否知道，还取决于它能否保存理由、分配任务、回应反证并修订公开结论，而非只看最后票数。

\minitopic{认识不正义损害共同证据库}证言不正义使某些说话者因身份偏见获得不当可信度亏欠；诠释不正义使共同概念资源不足以表达一类经验\parencite{fricker2007}。这不仅伤害个人尊严，也降低制度发现真相的能力：重要观察无法进入记录，异常被当作不专业，后来调查又因缺少早期材料而失败。多样性只有在成员拥有实际发言机会、带来独立相关信息、批评被处理并影响结论时才具有认识价值；象征性席位不会自动改善结果。

\minitopic{制度客观性来自可组织的批评}公开标准、原始数据版本、独立先行判断、受保护异议和事后校准把个人局限转化为可纠正过程\parencite{longino1990}。制度不可能消除所有偏见，却能决定偏见是否留下痕迹、是否被另一视角发现。好的程序还必须指定责任：谁保存来源，谁判断升级，谁有权暂停决定，何种新证据触发重开。没有责任人和触发条件的“欢迎批评”只是一种姿态。

\minitopic{阶段性结论}社会知识的核心问题不是个人知识加上社交背景，而是认识资格如何在说话者、听者、工具和制度之间生成与转移。可靠依赖要求初始信任、后来校准、来源独立性、专家元证据和可响应异议共同工作。真正的认识自主不是孤立，而是在不可避免的依赖中保留追踪、质疑和修正能力。
